\documentclass[11pt,a4paper]{article}
\usepackage[margin=2.2cm]{geometry}
\usepackage{amsmath,amssymb}
\usepackage{graphicx}
\usepackage{booktabs,tabularx,longtable}
\usepackage{xcolor}
\usepackage{float}
\usepackage{caption}
\usepackage{listings}
\usepackage{tikz}
\usetikzlibrary{arrows.meta,positioning,shapes.geometric,fit}
\usepackage[colorlinks=true,linkcolor=blue!50!black,urlcolor=blue!50!black,citecolor=blue!50!black]{hyperref}

\lstset{basicstyle=\ttfamily\small,breaklines=true,columns=fullflexible,frame=single,
        backgroundcolor=\color{gray!6},keepspaces=true,showstringspaces=false}
\usepackage{xurl}
\urlstyle{tt}
\newcommand{\code}[1]{\nolinkurl{#1}}
\usepackage{array}
\newcolumntype{L}[1]{>{\raggedright\arraybackslash}p{#1}}
\newcommand{\SK}{\textsc{sk}}
\newcommand{\fig}[3][0.98]{\begin{figure}[H]\centering\includegraphics[width=#1\linewidth]{#2}\caption{#3}\end{figure}}
\newcommand{\term}[1]{\textbf{#1}}

\graphicspath{{../SK_game/figures/}{../optimal_algorithms/figures/}{../embedding_explanation/figures/}{../noise_and_chains/figures/}}

\title{The SK game, its optimisation problem and quantum annealing on D-Wave\\[4pt]
\large What the project does, how every algorithm works, what noise does, and how it was tested}
\author{Project notes for \texttt{Projects\_master\_HPQC/scripts/QuantumComputing/Dimod}}
\date{}

\begin{document}
\maketitle
\tableofcontents
\newpage

%======================================================================================
\section{What this project is}
%======================================================================================

The starting point is the model of Garnier-Brun, Benzaquen and Bouchaud, \emph{Unlearnable games and
``satisficing'' decisions} (arXiv:2312.12252): $N$ agents take binary decisions, every agent's reward
depends on all the others through random couplings, and the agents \emph{learn} by remembering past rewards.
Inside the game sits an optimisation problem: the configuration of decisions that maximises the average
reward. That problem is the Sherrington--Kirkpatrick (\SK) spin glass, the textbook example of a hard
optimisation landscape.

The project does three things.
\begin{enumerate}
\item It \textbf{reproduces the game} (package \code{skgame}, all figures of the paper).
\item It \textbf{benchmarks optimisers} on the planner's problem: exact enumeration, greedy, the game's own learning,
simulated annealing, parallel tempering, extremal optimisation, tabu search, SDP relaxation, and
\textbf{quantum annealing on D-Wave}. It does this for growing $N$ and for parameter \emph{regimes} where the
optimum is easy or hard to retrieve.
\item It \textbf{runs this campaign} on a laptop or on Cineca's Leonardo, with the QPU reached through D-Wave Leap.
It compares the two D-Wave topologies, Zephyr (Advantage2) and Pegasus (Advantage), with an all-to-all embedding.
\end{enumerate}

\begin{table}[H]\centering\small
\begin{tabularx}{\linewidth}{L{5.2cm}X}
\toprule
where & what \\
\midrule
\code{unlearnable_games_explanation/skgame} & the model, theory, solvers (\code{optim/solvers.py}), embeddings
(\code{optim/embedding.py}), annealing back-ends (\code{optim/backends.py}), regimes (\code{optim/regimes.py}) \\
\code{unlearnable_games_explanation/main_simulation.py} & the staged campaign: \code{init}, \code{instances}, \code{classical},
\code{embed}, \code{qpu}/\code{sim}, \code{analyze} (+ notebook) \\
\code{unlearnable_games_explanation/local/} & local environment and one-command pipeline (\code{README_LOCAL.md}) \\
\code{unlearnable_games_explanation/cineca/} & Leonardo: offline environment, Slurm jobs, sync with a local machine (\code{README_CINECA.md}) \\
\code{visualization/SK_game} & the model step by step, glossary \\
\code{visualization/optimal_algorithms} & every algorithm, with traces of spin state and environment \\
\code{visualization/embedding_explanation} & how $K_N$ is embedded on Zephyr and Pegasus; chains and tuning \\
\code{visualization/noise_and_chains} & chain strength derivation, noise experiments \\
\bottomrule
\end{tabularx}
\caption{Map of the project. Every figure in this document is regenerated by the script of its folder.}
\end{table}

%======================================================================================
\section{The SK game, step by step}
%======================================================================================

\subsection{Agents, decisions, couplings}
Each agent $i=1,\dots,N$ takes a \term{decision} $S_i=\pm1$, its \term{spin}. The reward of agent $i$ is
$R_i=\sum_jJ_{ij}S_j$. The couplings are random,
\begin{equation}
J=\Big(1-\frac{\varepsilon}{2}\Big)J^S+\frac{\varepsilon}{2}J^A+\frac{J_0}{N},\qquad
J^S=J^{S\top},\ J^A=-J^{A\top},\ \text{entries}\sim\mathcal N(0,1/N).
\end{equation}
The \term{non-reciprocity} $\varepsilon$ controls how much $J_{ji}$ resembles $J_{ij}$; their correlation is
$\eta=(1-\varepsilon)/(1-\varepsilon+\varepsilon^2/2)$, equal to $1$, $0$, $-1$ at $\varepsilon=0,1,2$.
$J_0$ is a cooperative bias.

\fig{step01_couplings.png}{Couplings for $\varepsilon=0,1,2$ and the measured correlation $\eta$ between $J_{ij}$ and $J_{ji}$.}

\subsection{The network and frustration}
The planner only sees the symmetric part $W_{ij}=\frac12(J_{ij}+J_{ji})$ (zero diagonal). $W_{ij}>0$: $i$ and
$j$ prefer to agree; $W_{ij}<0$: they prefer to disagree. A loop of couplings whose signs cannot all be
satisfied is \term{frustrated}. In \SK\ about half of all triangles are frustrated, and frustration is the
root of the difficulty.

\fig{step02_network.png}{Eight agents. In the dashed triangle no choice satisfies all three pairs: at most 2 of 3.}

\subsection{Local field, satisfaction, energy, reward}
For a configuration $S$:
\begin{align}
h_i&=\sum_jW_{ij}S_j && \text{\term{local field}: the pull of the others on }i,\\
\lambda_i&=S_ih_i && \text{\term{satisfaction}: }>0\text{ satisfied},\ <0\text{ unsatisfied},\\
H(S)&=-\sum_{i<j}W_{ij}S_iS_j=-\tfrac12\sum_i\lambda_i && \text{\term{energy}},\\
R_N(S)&=\tfrac1NS^\top JS=\tfrac1NS^\top WS=-\tfrac{2H}{N} && \text{\term{average reward}}.
\end{align}
\textbf{Maximising the average reward is minimising the energy}: the planner's problem is the \SK\ ground state.

\fig{step03_configuration.png}{Left: a random configuration with unsatisfied agents (red bars). Right: after
flipping unsatisfied agents until none is left: everyone is satisfied and the energy is lower.}

\subsection{The single-flip move}
Flipping agent $k$ changes only the terms that contain $k$:
\begin{equation}
\Delta H_k=2\sum_jW_{kj}S_kS_j=2\lambda_k,\qquad\text{then }h_i\leftarrow h_i+2W_{ik}S_k^{\rm new}\ \ (O(N)).
\end{equation}
\textbf{A flip lowers the energy if and only if the agent is unsatisfied.} Almost every classical algorithm
below is built from this move.
\fig[0.6]{step04_single_flip.png}{The predicted $2\lambda_k$ and the measured energy change coincide.}

\subsection{The landscape: local minima, basins, barriers}
A \term{local minimum} (metastable state) is a configuration where every agent is satisfied
($\lambda_i\ge0$). No single flip improves it, but it need not be the \term{ground state}.
\fig{step05_landscape.png}{All $2^{14}$ configurations of one instance. The local minima (red) can be close in
energy to the ground state but several flips away: separated by energy barriers.}
\fig[0.6]{step06_basins.png}{Basins of steepest descent: a random start ends in the ground state with probability 0.33 for $N=14$.}
\fig[0.6]{step07_metastable_count.png}{Exact count of local minima: $\propto e^{0.199N}$. The ground state's basin
becomes an exponentially small fraction of the space.}
\fig[0.6]{step08_ground_state.png}{The ground-state energy per spin tends to the Parisi value $-0.7632$, with corrections $\approx0.75N^{-2/3}$.}

\subsection{The game: learning, phases, satisficing}
Each agent keeps an \term{incentive} $Q_i$ and plays $S_i=\mathrm{sign}\,Q_i$ (at temperature $T=1/\beta$:
$P[S_i=\pm1]=\frac12(1\pm\tanh\beta Q_i)$), then
\begin{equation}
Q_i(t+1)=(1-\alpha)Q_i(t)+\alpha R_i(t).
\end{equation}
Depending on $(\alpha,\beta,\varepsilon,J_0)$ the long-time behaviour is a fixed point, a 2-cycle, chaos, or a
paramagnet. These are the regimes of the campaign.
\fig{step09_game_dynamics.png}{The five regimes used in the campaign. Top: decisions over time; bottom: reward vs the optimum.}

At a fixed point $Q_i=R_i$ and $S_i=\mathrm{sign}\,R_i$, so $S_iR_i\ge0$. At $\varepsilon=0$, where $R_i=h_i$,
\textbf{every fixed point is a local minimum}. The game satisfies everyone, but it does not optimise: this is
\emph{satisficing}.
\fig{step10_satisficing.png}{At the game's fixed points no agent is unsatisfied; the median distance from the optimum is 9.4\% at $N=64$.}

\begin{table}[H]\centering\footnotesize
\begin{tabular}{lllll}
\toprule
regime & $(\varepsilon,J_0,\alpha,\beta)$ & planner problem & game phase & expected \\
\midrule
\code{ferro_easy} & $(0,3,0.1,\infty)$ & ferromagnetic & fixed point & easy for everybody \\
\code{sg_boundary} & $(0,1,0.1,\infty)$ & SG/ferro boundary & fixed point & intermediate \\
\code{sg_satisficing} & $(0,0,0.1,\infty)$ & spin glass & fixed point & hard; game below optimum \\
\code{sg_cycles} & $(0,0,1,\infty)$ & spin glass & 2-cycle & game never optimal \\
\code{nonreciprocal_chaos} & $(1,0,0.1,\infty)$ & spin glass & chaos & game never optimal \\
\code{paramagnet} & $(0,0,0.1,0.25)$ & spin glass & paramagnet & random play \\
\bottomrule
\end{tabular}
\caption{Regimes. The static difficulty depends only on $J_0^{\rm eff}=J_0/(1-\varepsilon/2)$: spin glass below 1, ferromagnet above.}
\end{table}

%======================================================================================
\section{The algorithms, in the language of the game}
%======================================================================================

Every algorithm answers one question: \emph{how do we get out of a local minimum?} For each one we give what it
does to the spins, how it escapes, its cost, where it is in the code, and two kinds of pictures:
\begin{itemize}
\item the \textbf{spin state and environment} figure (prefix \code{se_}). Three rows share the time axis: the
spin state (red = the agent agrees with the ground state $S^\star$, one global sign per run), the
\emph{environment} the algorithm reacts to, and the energy;
\item the \textbf{behaviour} figure (prefix \code{alg}): statistics over many runs.
\end{itemize}
The traces come from \code{optimal_algorithms/traced.py}, which has the same update rules as the benchmark
solvers. Steepest descent reproduces the package's \code{_greedy} exactly; the others reach the exact optimum
on small instances.

\subsection{Exact enumeration (Gray code)}
It visits all $2^{N-1}$ configurations ($S_{N-1}$ is fixed by the global flip symmetry) in an order where
consecutive configurations differ by one flip, so each step is a single-flip move. The cost is $N2^{N-1}$
operations: $24\cdot2^{23}\approx2.0\times10^8$ for $N=24$ (sub-second), $64\cdot2^{63}\approx5.9\times10^{20}$
for $N=64$ (impossible). This is why the exact reference stops at $N=24$.
\fig{alg01_exact.png}{Left: one spin changes per step. Right: all $2^{15}$ energies for $N=16$, and the best so far.}

\subsection{Greedy: steepest descent with restarts}
It flips the most unsatisfied agent until none is left, then restarts from a new random configuration. It
escapes only by restarting: every descent ends in the basin it started in.
\fig{se_greedy.png}{Spin state, satisfactions (environment) and energy during one descent. The final state is
satisfied everywhere and close in energy to $S^\star$, but many agents (blue) still disagree with it.}
\fig{alg02_greedy.png}{The probability that a descent ends in the ground state falls from 0.52 ($N=10$) to 0.15
($N=24$). The restarts needed for 99\% confidence, $r=\ln0.01/\ln(1-p)$, grow from 6 to 28.}

\subsection{The game's learning as an optimiser}
All agents update $Q_i\leftarrow(1-\alpha)Q_i+\alpha h_i$ together and play $\mathrm{sign}\,Q_i$. An agent changes
decision once it has been unsatisfied long enough for $Q_i$ to change sign. The run freezes in the first local
minimum: a local search with memory.
\fig{se_game.png}{The environment is the memory $Q_i$. Decisions stop changing after about 15 steps. The final
energy is within 1.4\% of the ground state, yet 11 of 24 agents disagree with $S^\star$: close in energy, far in configuration.}
\fig{alg03_game.png}{One game run finds the optimum with probability 0.11, one greedy descent 0.12. Greedy
polishing of the game's fixed point changes nothing, because it is already a local minimum.}

\subsection{Simulated annealing}
Metropolis single flips: an unsatisfied agent always flips, a satisfied one with probability
$e^{-\beta\Delta H}=e^{-2\beta\lambda_k}$. The inverse temperature $\beta$ rises from 0.1 to 5. When hot it crosses
barriers; as it cools, the barriers close one after another. The solver returns the best state seen, polished by greedy.
\fig{se_sa.png}{Environment: $\beta(t)$ and the satisfactions. Hot: a random mosaic of decisions. Cooling: groups of agents lock in and the satisfactions turn green.}
\fig{alg04_sa.png}{Acceptance falls from 0.95 to 0.05 with $\beta$. More sweeps give the ground state more often (0.19 at 25 sweeps, 0.98 at 800). The last state alone is a poor estimate (dashed).}

\subsection{Parallel tempering}
$R$ replicas at temperatures $\beta_1<\dots<\beta_R$. After each Metropolis sweep, neighbouring temperatures swap
configurations with probability $\min\{1,e^{(\beta_k-\beta_{k+1})(E_k-E_{k+1})}\}$. A trapped replica gets heated,
melts out of its minimum and cools into another one. Parallel tempering is also the reference solver above $N=24$.
\fig{se_pt.png}{Spin state of the coldest replica, and the environment: overlap with $S^\star$ of the replica at
each temperature. Good configurations appear at cold temperatures and are passed down the ladder.}
\fig{alg05_pt.png}{Energies at 12 temperatures; replicas random-walk across the ladder; swap rates 0.6--0.9.}

\subsection{$\tau$-extremal optimisation}
Rank the agents by satisfaction (rank 1 = least satisfied). Pick rank $k$ with probability $\propto k^{-\tau}$ and
flip it \emph{unconditionally}. There is no temperature and no acceptance test, so the state never freezes; the best visit is kept.
\fig{se_eo.png}{Black dots: the agent flipped at each step, usually among the least satisfied.}
\fig{alg06_eo.png}{Chosen ranks follow $k^{-1.3}$. The best $\tau$ is around 1.2--1.6.}

\subsection{Tabu search}
At every step flip the allowed agent with the smallest $\Delta H=2\lambda_k$, even uphill. A flipped agent
becomes \emph{tabu} for \code{tenure} steps, so the search cannot fall back into the minimum it just left.
D-Wave's \code{TabuSampler} runs the MST2 multistart variant with tenure $\min(20,\lfloor N/4\rfloor)$.
\fig{se_tabu.png}{Black dots: agents that are tabu. The energy keeps oscillating close to the ground state, which it visits repeatedly.}
\fig{alg07_tabu.png}{Without memory (tenure 0) it finds the ground state in 20\% of runs; with tenure 4--8, in 97--100\%.}

\subsection{SDP relaxation and Goemans--Williamson rounding}
Replace $S_i$ by unit vectors $v_i\in\mathbb R^k$ and minimise $-\frac12\sum W_{ij}v_i\cdot v_j$, which is convex in
$X=VV^\top$. The mixing method aligns each vector with its vector local field, $v_i\leftarrow$
normalise$(\sum_jW_{ij}v_j)$. Rounding with a random hyperplane gives $S_i=\mathrm{sign}(v_i\cdot r)$. The relaxed
minimum is a \textbf{certified lower bound} on $H_{\min}$.
\fig{se_sdp.png}{Environment: alignment of each vector with the ground-state direction. The relaxed energy settles below the ground state (the bound).}
\fig{alg08_sdp.png}{The bound $-0.64N$ lies below $H_{\min}=-0.557N$. Rounding alone hits the optimum 3\% of the time; with greedy polishing, 17\%.}

\subsection{The dimod samplers}
\code{dwave.samplers} implements SA (C++, sweeps in fixed order, $\beta$ range computed from the instance),
tabu and steepest descent behind the same \code{dimod} interface as the QPU.
\fig{alg09_dwave.png}{$N=64$: SA finds the best known energy in 5/54/93\% of reads with 10/100/1000 sweeps; tabu in every read.}
\fig[0.8]{alg11_overview.png}{All classical solvers on the same instances (quick budgets). PT, dwave SA and tabu stay at 1 up to $N=64$; greedy and the game's learning degrade first.}

\subsection{Quantum annealing}
The spins become qubits. The machine starts in the ground state of a transverse field, where every configuration
is superposed, and slowly turns it off while turning on the \SK\ energy:
\begin{equation}
\mathcal H(s)=-A(s)\sum_i\sigma^x_i+B(s)H_P(\sigma^z),\qquad s=t/T:0\to1 .
\end{equation}
Measuring at $s=1$ gives one configuration (a \emph{read}). If $T\gg1/\Delta_{\min}^2$, where $\Delta_{\min}$ is
the minimum gap of $\mathcal H(s)$, the system stays in its ground state (adiabatic theorem). Instead of climbing
barriers, it \emph{tunnels} through them.
\fig{se_qa.png}{Exact Schr\"odinger simulation, 8 qubits, idealised schedule $A=1-s$, $B=s$. Spin state: the
expected agreement of each agent with $S^\star$ (0 = undecided). The agents start undecided and decide together.
Probability flows from all levels to the lowest ones. With $T=100<1/\Delta_{\min}^2\approx700$, two agents often
end on the wrong side and $P(\text{ground})=0.35$.}
\fig{alg10_qa.png}{Spectrum (minimum gap 0.038), success vs anneal time (0.29 at $T=10$, 0.89 at $T=1000$), and the distribution of read energies.}

%======================================================================================
\section{How D-Wave is used in this project}
%======================================================================================

\subsection{The path of one QPU call}
\begin{figure}[H]\centering
\begin{tikzpicture}[node distance=0.55cm and 0.5cm,
  box/.style={draw,rounded corners,fill=blue!6,align=center,font=\small,minimum height=1.05cm,text width=2.6cm},
  arr/.style={-{Stealth},thick}]
\node[box] (p) {\code{PlannerProblem}\\$W_{ij}$, $N$};
\node[box,right=of p] (b) {\code{to_bqm()}\\dimod BQM (SPIN)\\$E=H(S)$};
\node[box,right=of b] (e) {embedding\\\code{CliqueEmbedding}\\chains $C_i$};
\node[box,right=of e] (c) {\code{embed_bqm}\\split $W_{ij}$, chains $-c$};
\node[box,below=of c] (s) {\code{auto_scale}\\divide by $\max|J|\approx c$};
\node[box,left=of s] (q) {QPU anneal\\\code{num_reads},\\\code{annealing_time}};
\node[box,left=of q] (u) {unembed\\\code{chain_break_method}};
\node[box,left=of u] (st) {\code{sampleset_statistics}\\rewards, counts,\\timing, breaks, polish};
\draw[arr] (p)--(b); \draw[arr] (b)--(e); \draw[arr] (e)--(c); \draw[arr] (c)--(s);
\draw[arr] (s)--(q); \draw[arr] (q)--(u); \draw[arr] (u)--(st);
\end{tikzpicture}
\caption{One call of \code{AnnealBackend.sample}. The embedding and scaling happen inside
\code{FixedEmbeddingComposite} and \code{DWaveSampler}; unembedding and chain-break resolution happen on the way back.}
\end{figure}

\subsection{Account and access: Leap}
The QPUs are in D-Wave's cloud, \emph{Leap}. Cineca grants QPU time through a Leap project; access needs only
an API \emph{token} for that project and outbound HTTPS, from any machine (Leonardo login node, the serial
partition, or a laptop). Command-line tools from \code{dwave-cloud-client}:
\begin{lstlisting}
dwave config create     # store the token in ~/.config/dwave/dwave.conf
dwave solvers -la       # list QPUs: Advantage_* (Pegasus), Advantage2_* (Zephyr)
dwave ping              # submit a tiny problem (a few ms of QPU time)
\end{lstlisting}

\paragraph{Getting the token, step by step} (details in \code{Dimod/LEAP_TOKEN.md}). A free Leap account has
no usable token: the token belongs to a Leap \emph{project} with QPU time, which Cineca grants.
\begin{enumerate}
\item Create a Leap account at \url{https://cloud.dwavesys.com/leap/} with your \textbf{institutional} e-mail; every user needs one.
\item Apply to Cineca's \textbf{ISCRA-C} call for D-Wave time. List the institutional e-mails of the Leap accounts and the minutes per month (Cineca's page states at most 20 min/month for 3 months, evaluated monthly). The \code{full} campaign needs about $124\,$s $\approx2\,$min, so 5--10 min/month is ample.
\item After approval your account joins a Leap project with the minutes. Select it in Leap (username menu $\to$ Projects): each project has its own token.
\item Copy the token from the dashboard (left column, \emph{API Token}, COPY), or from a terminal:
\code{dwave auth login --oob} (open the printed link on your laptop, paste back the code), \code{dwave leap project ls}, then \code{dwave leap project token --project <NAME>}.
\item Store it on every machine that calls the QPU: \code{dwave setup --oob --project <NAME>}, or \code{dwave config create} (paste it), or \code{export DWAVE_API_TOKEN=...}. Check with \code{dwave solvers -la} and \code{main_simulation.py qpu --check}.
\end{enumerate}
Never commit the token. If it leaks, RESET it on the Leap dashboard.

\subsection{Functions and parameters, and their role here}
{\small\begin{longtable}{L{4.0cm}L{5.4cm}L{5.6cm}}
\toprule
function / object & what it does & parameters and settings used here \\
\midrule\endhead
\code{dimod.BinaryQuadraticModel} (via \code{PlannerProblem.to_bqm}) & the problem in Ising form: linear $h$, quadratic $J$, offset; SPIN variables & $h=0$, $J_{ij}=-W_{ij}$ ($i<j$), so the BQM energy is exactly $H(S)$ \\
\code{dwave.system.DWaveSampler} & connection to one QPU solver; exposes \code{properties}, \code{parameters}, \code{to_networkx_graph()} (working qubits and couplers) & \code{solver={'topology__type': 'zephyr'}} or \code{'pegasus'} chooses the family; \code{--solver NAME} pins a specific QPU \\
\code{to_networkx_graph()} & the real hardware graph, with dead qubits removed & the embedding is computed on it, so it fits the chip used \\
\code{minorminer.busclique.find_clique_embedding} & native all-to-all embedding: chains = one vertical + one horizontal line segment & \code{use_cache=True}; result stored in \code{embeddings/<solver>/clique/N....json} \\
\code{busgraph_cache(G).largest_clique()} & largest $K_N$ the chip supports & larger $N$ are skipped automatically; printed by \code{qpu --check} \\
\code{minorminer.find_embedding} & general heuristic embedding & baseline only (\code{qpu:zephyr:minorminer}) \\
\code{dwave.system.FixedEmbeddingComposite} & embeds the BQM with a given embedding, samples, unembeds & used for both topologies, so the two runs differ only in hardware \\
\code{chain_strength} & strength $c$ of the couplers inside chains & \code{partial(uniform_torque_compensation, prefactor=p)}, $c=p\,J_{\rm rms}\sqrt{\bar d}\approx p$ for \SK; $p$ = \code{chain_strength_prefactor}, default 1.414 (\S\ref{sec:chains}) \\
\code{chain_break_method} & how a broken chain is read & \code{majority_vote} (default), \code{minimize_energy} (wrapped in a named function, see \S\ref{sec:testing}), \code{weighted_random}, \code{discard} \\
\code{return_embedding=True} & puts the embedding context (incl.\ the chain strength used) in \code{sampleset.info} & recorded in every result file \\
\code{num_reads} & anneal--readout cycles per call & 1000 (\code{full}); sets the statistics of $p_{\rm success}$ per read and the QPU time \\
\code{annealing_time} & duration of one anneal, $\mu$s & 20; longer is closer to adiabatic, but costs more per read \\
\code{auto_scale} (default on) & divides $h,J$ to fit the programmable range $[-1,1]$ & the reason for the compression in \S\ref{sec:chains} \\
\code{SampleSet} & reads: \code{record.sample}, \code{record.energy}, \code{num_occurrences}, \code{chain_break_fraction}, \code{info['timing']} & \code{aggregate()} merges identical reads; rewards recomputed from the spins \\
\code{info['timing']} & \code{qpu_access_time}, \code{qpu_programming_time}, \code{qpu_sampling_time}, per-sample anneal/readout/delay ($\mu$s) & \code{qpu_access_time} is what Leap charges; logged per call \\
\code{solver.estimate_qpu_access_time} & D-Wave's own timing model & used by \code{qpu --check}, spends nothing \\
\code{dwave.embedding.embed_bqm} & the embedding step itself & used in \code{noise_and_chains} to measure the programmed couplings \\
\code{dwave.samplers.SimulatedAnnealingSampler}, \code{TabuSampler}, \code{SteepestDescentSolver} & classical samplers with the dimod interface & classical baselines; \code{sim:*} back-ends \\
\code{dimod.StructureComposite} & restricts a sampler to a hardware graph & \code{sim:zephyr:clique}: the same embedded problem solved by SA, no token \\
\code{dwave.system.testing.MockDWaveSampler} & offline stand-in for a QPU with a given topology & tests of the whole QPU path \\
\code{dwave.preprocessing.SpinReversalTransformComposite} & gauge transforms & mitigation of systematic biases (\S\ref{sec:noise}); not yet wired into the back-ends \\
\bottomrule
\end{longtable}}

\subsection{The back-ends of the campaign}
A back-end is a string \code{device:topology:strategy}:
\begin{itemize}
\item \code{qpu:zephyr:clique}: Advantage2, Zephyr, degree 20, clique chains $\approx N/16+2$;
\item \code{qpu:pegasus:clique}: Advantage, Pegasus, degree 15, chains $\approx N/12+2$;
\item \code{sim:*}: classical stand-ins.
\end{itemize}
Both QPU back-ends solve the same instances with the same \code{AnnealSettings}. \code{analyze} writes
\code{backend_comparison.csv} with the winner per (problem class, $N$): higher success per read, ties broken by
lower residual.

\subsection{Why the QPU budget doubles with two topologies}
A QPU call is one problem on one machine. The comparison needs \emph{every} instance solved on \emph{both} machines:
\begin{equation}
n_{\rm call}^{\rm per\ back\text{-}end}=|N_s|\cdot n_{\rm classes}\cdot n_{\rm inst}=20\cdot4\cdot5=400 .
\end{equation}
Nothing can be shared between the two back-ends. Zephyr and Pegasus are different QPUs, the embeddings
differ, and each call reprograms its own chip. The time of one call is programming plus one anneal--readout--delay cycle per read:
\begin{equation}
t_{\rm call}\approx t_{\rm prog}+n_{\rm reads}(t_a+t_{\rm ro}+t_d)=15\,{\rm ms}+1000\cdot(20+100+20)\,\mu{\rm s}
=15\,{\rm ms}+140\,{\rm ms}=155\,{\rm ms}.
\end{equation}
Hence
\begin{equation}
T=2\cdot400\cdot0.155\,{\rm s}=2\cdot62\,{\rm s}=124\,{\rm s}.
\end{equation}
These are typical values: readout dominates and differs between systems, and \code{qpu --check} gives each
solver's own estimate. To spend less, keep the comparison fair and cut what both share: fewer reads (time is
almost linear in \code{num_reads}), fewer instances, a coarser $N$ grid, or one topology only
(\code{init --backends qpu:zephyr:clique}).

%======================================================================================
\section{Embedding: all-to-all on Zephyr and Pegasus}
%======================================================================================

\SK\ couples every pair of spins ($K_N$), but each qubit has only 15 (Pegasus) or 20 (Zephyr) couplers. Each
spin becomes a \term{chain} of qubits.
\fig[0.5]{step01_logical_problem.png}{The problem graph: $N(N-1)/2$ couplings.}
\fig{step02_hardware.png}{Qubits of two orientations (vertical, horizontal).}
\fig{step03_couplers.png}{Three kinds of coupler: internal (vertical $\times$ horizontal), external (along a line), odd. Degree $16+2+2=20$ on Zephyr, $12+2+1=15$ on Pegasus.}
\fig{step04_lines.png}{Qubits joined end to end by external couplers form a line.}
\fig{step05_one_chain.png}{A chain is one vertical line segment plus one horizontal segment, joined by a single internal coupler (checked for all 160 chains at $N=160$ on both topologies).}
\fig{step06_all_chains.png}{The full clique embedding on a small graph.}
\fig{step07_crossbar.png}{Why it is all-to-all: on one side of the staircase of corners, every vertical segment meets every horizontal one, so each pair $\{i,j\}$ has a $V_i$--$H_j$ or $V_j$--$H_i$ coupler.}
\fig{step08_link_matrix.png}{Couplers between every pair of chains, $K_{64}$: all off-diagonal entries are $\ge1$.}
\fig{step09_scaling.png}{Chain length and qubits vs $N$. At $N=160$: Zephyr 12, Pegasus 15. The heuristic embedder gives longer, uneven chains.}
\fig[0.6]{step10_defects.png}{Missing qubits shrink the largest clique (here with random defects); Zephyr keeps its lead.}

\paragraph{The algorithm (busclique).}
\begin{enumerate}
\item Group qubits into lines.
\item Keep only intact segments (all qubits and external couplers working).
\item In a block of the chip, pair $N$ vertical segments with $N$ horizontal ones along a staircase of corners.
\item The chains are $C_i=V_i\cup H_i$.
\item Optimise the placement by dynamic programming, for the largest $N$ and the shortest chains.
\item Cache the result.
\end{enumerate}
The heuristic \code{find_embedding} (Cai--Macready--Roy) works differently. It routes each chain as a tree, temporarily allowing shared qubits with an exponential penalty, and keeps re-routing until no qubit is shared.

%======================================================================================
\section{Chains and chain strength: how it was worked out}\label{sec:chains}
%======================================================================================

This section reconstructs the reasoning in the order it actually happened. Every number is reproduced by \code{noise_and_chains.py}.

\paragraph{1. The question.} Inside a chain the couplers are set to $-c$. How large is $c$ by default, and what does that do to the rest of the problem?

\paragraph{2. Read the function.} \code{inspect.getsource(uniform_torque_compensation)} shows that the default is
\begin{equation}
c=p\cdot J_{\rm rms}\cdot\sqrt{\bar d},\qquad J_{\rm rms}=\sqrt{\tfrac{1}{|E|}\textstyle\sum J_{ij}^2},\ \ \bar d=\text{average degree},\ \ p=1.414 .
\end{equation}
The docstring's argument: half a chain feels a random torque from its neighbours, a sum of about $\bar d/2$ couplings of random sign, of size $\propto J_{\rm rms}\sqrt{\bar d}$. The chain coupling must exceed it.

\paragraph{3. Evaluate it for \SK.} $J_{\rm rms}=1/\sqrt N$ (entries $\mathcal N(0,1/N)$) and $\bar d=N-1$ (complete graph), so
\begin{equation}
c=p\frac{1}{\sqrt N}\sqrt{N-1}=p\sqrt{1-1/N}\to p\approx1.4 .
\end{equation}
Checked against the library on four problem classes. The script asserts that my formula and the library agree.

\paragraph{4. Meaning.} $c/p=J_{\rm rms}\sqrt{N-1}$ is the standard deviation of a local field $h_i$ (measured 0.991 at $N=160$). \textbf{The chain strength is $p$ times the typical pull of the other agents on one agent.}
\fig{A1_field_and_chain_strength.png}{$J_{\rm rms}\sqrt N\approx1$; $c$ is flat in $N$ ($\approx0.7$ at $\varepsilon=1$, where $W$ is halved); the local-field distribution with $\pm c$.}

\paragraph{5. Embedding splits the couplings.} $W_{ij}$ is divided evenly over the $|E_{ij}|$ couplers between two chains. On average $|E_{ij}|=1.21$ (Zephyr) and $1.11$ (Pegasus) at $N=160$, which reduces the rms by a factor 0.895 on Zephyr.

\paragraph{6. Auto-scaling divides by $c$.} The chain couplers ($1.398$) are much larger than any logical one ($\max|W|=0.29$), so $\lambda=1/c=0.715$.

\paragraph{7. Multiply.}
\begin{align}
\text{Zephyr, }N=160:&\quad 0.0784\times0.895\times0.715=0.0702\times0.715=0.0502,\\
\text{in general:}&\quad \frac{J_{\rm rms}}{c}=\frac{1}{p\sqrt{N-1}}=\frac{1}{1.414\cdot12.61}=0.056\quad(N=160).
\end{align}
\fig[0.65]{A2_compression_chain.png}{The programmed rms of the \SK\ couplings vs $N$ on both topologies, against $1/(p\sqrt{N-1})$.}
\fig{A3_programmed_couplings.png}{What the QPU receives: every chain coupler at $-1$, the whole instance in a band of rms 0.05 around 0.}

%======================================================================================
\section{Noise: what it does, what problems it raises, how to reduce it}\label{sec:noise}
%======================================================================================

\textbf{Important.} No QPU has been run in this project yet: there was no token in the development environment.
The experiments below are \emph{simulations}. Each one isolates a single mechanism with a deliberately simple
model, so they show the \emph{kind} of problem and the \emph{direction} of each remedy, not the QPU's numbers.
What I learned from them is stated with that limitation.

\subsection{Random errors on the couplings}
\textbf{Model.} Each coupler is off by $\delta\sim\mathcal N(0,\sigma)$ in QPU units, i.e.\ $\sigma c$ on every $W_{ij}$. The noisy problem is then solved exactly.
\fig{B1_coupling_noise.png}{At $\sigma=0.03$ the true optimum survives in 77\%, 60\%, 37\% of instances for $N=10,14,18$. Even a perfect solver of the noisy problem returns a configuration 1--2\% above the minimum.}
\textbf{What it produces.} The device solves a \emph{different} problem.
\textbf{The problem it raises.} The error is invisible from inside: reads are consistent with each other, and repeating the call does not help.
\textbf{Why it grows with $N$.} Relative to a typical coupling the error is
\begin{equation}
\frac{\sigma c}{J_{\rm rms}}=\sigma p\sqrt{N-1};\qquad N=160,\ \sigma=0.02:\ 0.02\cdot1.414\cdot12.61=0.36 .
\end{equation}
\fig[0.6]{B2_relative_noise_vs_N.png}{The compression forced by the chains turns a fixed analogue error into a relative error $\propto\sqrt N$.}

\subsection{Chain breaks, and chains that are too rigid}
\textbf{Model.} The embedded problem on ideal Zephyr, solved by classical SA as a stand-in, scanning $p$.
\fig{B3_chain_breaks.png}{Breaks: 25\% of chains at $p=0.25$, below 1\% for $p\ge1$. The raw success peaks at intermediate $p$ and falls to 0 at $p=3$, because rigid chains freeze single flips. Greedy polishing lifts the success from $\le0.05$ to $\approx0.2$.}
\textbf{Problem.} Two opposite failure modes, so the best $p$ depends on the instance.
\textbf{Remedies.} Scan $p\in[0.5,2]$; use \code{minimize_energy} where chains break; polish the reads; prefer shorter chains (Zephyr).

\subsection{Finite effective temperature}
\textbf{Model.} Reads drawn from $P(S)\propto e^{-\beta_{\rm eff}H(S)/c}$. The $1/c$ is the auto-scaling, which raises the effective temperature by $c$.
\fig{B4_temperature.png}{Hot reads: success falls from 0.27 to 0.11 between $N=8$ and 18. Per instance, some gaps are almost zero: there any temperature splits the reads.}
\textbf{Remedies.} Longer \code{annealing_time}; many reads (judge by TTS$_{99}$); polishing.

\subsection{Systematic biases and spin-reversal transforms}
\textbf{Model.} A fixed field offset per qubit, the same in every call. A gauge $g\in\{\pm1\}^N$ programs $W_{ij}g_ig_j$ and maps back $S_i=g_iS'_i$: the problem is unchanged, but the offset changes sign pattern.
\fig[0.65]{B5_bias_and_gauges.png}{Bias $\sigma=0.2$: the optimum is found with one gauge in 52\% of instances, with 2/4/8/16 gauges in 72/88/95/100\%.}

\subsection{Summary}
\begin{table}[H]\centering\small
\begin{tabularx}{\linewidth}{p{3.1cm}XX}
\toprule
noise source & symptom & remedy (knob) \\
\midrule
coupler errors & a reproducible but suboptimal best read; worse as $\sqrt N$ & shorter chains (Zephyr); smaller \code{chain_strength_prefactor} if chains allow it; polishing \\
chain breaks & \code{chain_break_fraction} above a few \% & scan the prefactor; \code{minimize_energy}; polishing \\
rigid chains & low success with zero breaks & lower the prefactor \\
temperature & success per read falls with $N$, varies between instances & longer \code{annealing_time}; more \code{num_reads}; polishing \\
systematic bias & results depend on the call & spin-reversal transforms, several calls \\
too-fast anneal & success grows with \code{annealing_time} & scan the anneal time, compare TTS$_{99}$ \\
\bottomrule
\end{tabularx}
\end{table}

%======================================================================================
\section{How everything was tested, and what testing taught}\label{sec:testing}
%======================================================================================

\subsection{The testing procedure}
\begin{enumerate}
\item \textbf{Unit and theory tests} (\code{pytest}, 107 tests). They cover the model (ensemble moments, limits, recorders, fixed points), the theory (13 symbolic checks, $T_c$, complexity vs brute force, DMFT vs simulation), the optimisers (exact is exact, heuristics reach the optimum on small instances, SDP bound, dimod wrappers), the output registry, and two new files:
\begin{itemize}
\item \code{test_embedding.py}: clique embeddings valid on both topologies; Zephyr chains shorter at production size; cache round trip; back-end spec parsing; the embedded simulator finds the exact ground state, with the chain strength equal to the formula;
\item \code{test_campaign.py}: regimes; full read statistics; the whole campaign end to end with resume; both QPU back-ends against \code{MockDWaveSampler} (D-Wave's offline mock QPU) with the comparison output.
\end{itemize}
\item \textbf{QPU path without a token.} \code{MockDWaveSampler} with Zephyr and Pegasus topologies goes through embedding on the mock's graph, chain strength, unembedding, timing fields and chain-break fraction.
\item \textbf{End-to-end runs.} The \code{smoke} and \code{small} presets with the simulated back-ends, the notebook executed headless, and \code{local/run_pipeline.sh} (20 s and 3 min).
\item \textbf{Environment.} The wheelhouse was built for Python 3.11 and 3.13 and installed with \code{--no-index} (nothing downloaded), after which the test suite passed. \code{local/setup_env.sh} was run on a clean copy.
\item \textbf{Figures as tests.} Every figure was looked at against the claim in its caption. Where they disagreed, the code or the claim was fixed (next subsection).
\end{enumerate}

\subsection{What testing found (experience)}
\begin{itemize}
\item \textbf{Seed sent to the QPU.} The original dimod wrapper always passed \code{seed}. The QPU rejects unknown parameters server-side, not with a Python \code{TypeError}, so the existing fallback would never have triggered. Fix: pass \code{seed} only to samplers that declare it.
\item \textbf{\code{MinimizeEnergy} has no \code{__name__}.} \code{unembed_sampleset} needs a named function. Found by the parametrised test; fixed with a wrapper.
\item \textbf{The simulator lost the topology metadata.} A graph rebuilt from the edges of \code{StructureComposite} lacks the \code{family} attribute that busclique needs. Fix: keep the ideal graph object itself.
\item \textbf{A claim that was too strong.} I first wrote that ``every vertical segment crosses every horizontal one''. The crossbar figure showed that this holds only on one side of the staircase of corners, which is enough for all-to-all. The text was corrected to what the data show.
\item \textbf{Measuring the wrong thing.} The first SA success curve was non-monotonic because it scored the \emph{last} state, while the solver returns the \emph{best} state seen, polished. Re-measured properly: monotonic, 0.19$\to$0.98.
\item \textbf{Too few samples.} The ground-state energy vs $N$ with 6 instances was dominated by sample-to-sample fluctuations. With 40 instances it follows the known finite-size law.
\item \textbf{Mixing two sources of randomness.} The first gauge experiment redrew the bias for every number of gauges, giving a non-monotonic curve. A systematic bias must be \emph{fixed}; with that, the expected monotonic gain appears.
\item \textbf{A confusing representation.} Fixing the global sign $S\to-S$ at every time step made whole columns flip colour. One sign per trajectory fixed it.
\end{itemize}
The general lesson: in this problem many plausible-looking curves are artefacts of what is measured (last vs best state, per-step vs per-run sign, redrawn vs fixed noise). Each figure needs a check against an independent prediction: an exact enumeration, a formula, or a library value.

%======================================================================================
\section{Running it}
%======================================================================================
\textbf{Locally} (\code{README_LOCAL.md}):
\begin{lstlisting}
bash local/setup_env.sh && source .venv/bin/activate
dwave config create                                  # only for the QPU
PRESET=small bash local/run_pipeline.sh              # ~3 min, simulated annealers
PRESET=full ANNEALER=both bash local/run_pipeline.sh # production, Zephyr + Pegasus
\end{lstlisting}
\textbf{On Leonardo} (\code{README_CINECA.md}): build the wheelhouse locally (\code{cineca/build_wheelhouse.sh 3.11}),
upload it, \code{cineca/install_offline.sh}, then \code{init}, \code{embed}, \code{qpu --check}, \code{cineca/submit_all.sh}.
The QPU stage can also run on the laptop: \code{cineca/sync.sh pull-inputs}, \code{qpu}, \code{push-anneal}.

\paragraph{Suggested first QPU session.}
\begin{enumerate}
\item \code{qpu --check}.
\item Two calls per topology (\code{QPU_CALLS=2}).
\item A chain-strength scan $p\in\{0.5,0.75,1,1.41,2\}$ $\times$ \{\code{majority_vote}, \code{minimize_energy}\} on 2--3 instances at the $N$ where success starts to drop.
\item An \code{annealing_time} scan.
\item The full campaign with the chosen settings.
\end{enumerate}

%======================================================================================
\appendix
\section{Glossary}
%======================================================================================
\begin{longtable}{p{4.2cm}p{10.8cm}}
\toprule
term & meaning \\ \midrule\endhead
spin, decision $S_i$ & the binary choice of agent $i$ \\
coupling $J_{ij}$, $W_{ij}$ & effect of $j$ on $i$'s reward; symmetric part seen by the planner \\
local field $h_i$ & $\sum_jW_{ij}S_j$, the pull of the others \\
satisfaction $\lambda_i$ & $S_ih_i$; a flip of $i$ changes the energy by $2\lambda_i$ \\
energy $H$, reward $R_N$ & $H=-\frac12\sum\lambda_i$, $R_N=-2H/N$ \\
local minimum & every $\lambda_i\ge0$; fixed point of the zero-temperature game at $\varepsilon=0$ \\
ground state & minimum of $H$; $H/N\to-0.7632$ \\
basin, barrier & starts that descend to a minimum; energy to climb between minima \\
frustration & loops of couplings that cannot all be satisfied \\
temperature, replica & acceptance of uphill moves; a copy of the system at one temperature \\
chain, chain strength $c$ & qubits representing one spin; coupling $-c$ holding them together \\
auto-scale & division of the programmed problem into $[-1,1]$ (by $\approx c$ here) \\
read, $p_{\rm success}$, TTS$_{99}$ & one anneal + measurement; fraction of reads at the optimum; $t_{\rm read}\ln0.01/\ln(1-p)$ \\
gauge (spin-reversal) & $S_i\to g_iS_i$ with $W_{ij}\to g_ig_jW_{ij}$: same problem, different hardware signs \\
\bottomrule
\end{longtable}

\end{document}
